\documentclass[11pt,titlepage,a4paper]{report}

%INCLUSIONE PACCHETTI
%---------------------------------------------
\usepackage[italian]{babel}
\usepackage{fancyhdr}
\usepackage{graphicx}
\graphicspath{{./pics/}}	% cartella di salvataggio immagini

% STILE DI PAGINA
%---------------------------------------------
\pagestyle{fancy}
\renewcommand{\sectionmark}[1]{\markright{\thesection.\ #1}}
\lhead{\nouppercase{\rightmark}}
\rhead{\nouppercase{\leftmark}}
\renewcommand{\chaptermark}[1]{%
\markboth{\thechapter.\ #1}{}}

%Ridefinisco lo stile plain della pagina
\fancypagestyle{plain}{%
	\lhead{\includegraphics[height=50pt]{logo.eps}}
	\chead{}
	\rhead{HappyCode inc \\ happycodeinc@gmail.com}
	\lfoot{BR-jsys}
	\cfoot{\thepage}
	\renewcommand{\headrulewidth}{1pt}
	\renewcommand{\footrulewidth}{1pt}
}
%---------------------------------------------


%INIZIO DOCUMENTO


%---------------------------------------------
% PRIMA PAGINA
%---------------------------------------------
\begin{document}

\hyphenation{re-qui-si-ti in-se-ri-men-to}

\begin{titlepage}
\begin{center}
\vspace*{0.5in}
\includegraphics{logo.eps}
\vspace*{0.2in}

{\Large \textbf{BR-jsys}}

{\Large \emph{Business Rules} per sistemi gestionali in architettura J2EE } 
\vspace{1.3in}
\par\rule{10cm}{.4pt} \par
\par\rule{12cm}{1pt} \par
\vspace*{0.5in}
\LARGE \textbf {PIANO DI QUALIFICA}
\vspace*{0.5in}
\par\rule{12cm}{1pt} \par
{\large Versione 0.2 - \today}
\par\rule{10cm}{.4pt} \par

\end{center}
\end{titlepage}
\vspace*{0.5in}

%---------------------------------------------
% SECONDA PAGINA
%---------------------------------------------
\begin{center}
\thispagestyle{plain}
\begin{table}[htbp]
\large{
\begin{tabular}{l}
\Large{\textbf{\textsf{Capitolato: ''BR-jsys``}}} \\
\begin{tabular}{||p{6cm}||p{6cm}||}
\hline
\textbf{Data creazione:} & 19/11/07 \\
\hline
\textbf{Versione:} & 0.2 \\
\hline
\textbf{Stato del documento:} & formale, esterno \\
\hline
\textbf{Redazione:} & Luca Appon \\
\hline
\textbf{Revisione:} & Marco Tessarotto   \\
\hline
\textbf{Approvazione:}  & Elena Trivellato\\
\hline
\end{tabular} \\
\end{tabular}
}
\end{table}

\begin{table}[hbtp]
\large{
\begin{tabular}{l}
\Large{\textbf{\textsf{Lista di distribuzione}}} \\
\begin{tabular}{||p{6cm}||p{6cm}||}
\hline
{Elena Trivellato}& Responsabile di progetto \\ 
\hline 
{Filippo Carraro}& Amministratore \\ 
\hline
{Marco Tessarotto}& Verificatore \\
\hline
{Tullio Vardanega}& Committente \\
\hline
{Renato Conte}& Committente \\
\hline
{Zucchetti S.r.l.}& Proponente \\
\hline
\end{tabular} \\
\end{tabular}
}
\end{table}

\begin{table}[hbtp]
\large{
\begin{tabular}{l}
\Large{\textbf{\textsf{Diario delle modifiche}}} \\
\begin{tabular}{||p{2cm}||p{3.5cm}||p{6cm}||}
\hline
\textbf{Versione} & \textbf{Data rilascio} & \textbf{Descrizione} \\
\hline
0.2 & 2007/12/06& Correzione errori. \\
\hline
0.1 & 2007/11/19 & Stesura preliminare del documento. \\
\hline
\end{tabular} \\
\end{tabular}

}
\end{table}
\end{center}

\tableofcontents	%Crea l'indice

\chapter{Introduzione}
\section{Scopo del documento}
Il presente documento identifica e spiega nel dettaglio le strategie di verifica e validazione proposte per il collaudo del sistema. Useremo l'analisi statica preliminare e l'analisi dinamica del codice per soddisfare i requisiti quali funzionalit\`a,affidabilit\`a e usabilit\`a del prodotto ''BR-jsys``.

\section{Scopo del prodotto}
Il prodotto richiesto verr\`a inserito nell'ambito di un progetto pi\`u ampio. Il suo scopo \`e quello di automatizzare il sistema di validazione dei dati in ingresso al database dell'applicazione principale.

\section{Glossario}
Viene fornito come documento esterno chiamato Glossario\_0\_4.pdf.

\chapter[Strategia di verifica]{Visione generale della strategia di verifica}

\section[Organizzazione, pianificazione, responsabilit\`a]{Organizzazione, pianificazione strategica e temporale, responsabilit\`a}
\subsection{Ciclo di vita}
Le attivit\`a seguiranno un modello di tipo evolutivo che permetter\`a di apportare eventuali modifiche ai vari documenti. L'obiettivo \`e comprendere le richieste del cliente e quindi sviluppare una migliore definizione dei requisiti del sistema. Il prototipo si concentra sulle parti che sono mal comprese con l'obiettivo di contribuire a chiarire i requisiti. In quest'ottica il prototipo si pu\`o trasformare progressivamente nel prodotto. L'obiettivo del processo di sviluppo \`e lavorare in stretto contatto con il cliente per indagare i requisiti e giungere ad un prodotto finale. Si svilupperanno quindi le parti del sistema che sono ben chiare (requisiti ben compresi) e, solo successivamente si aggiungeranno nuove parti/funzionalit\`a come proposto dal cliente. 
\subsection{Pianificazione delle attivit\`a}
In una prima fase gli analisti discuteranno e cercheranno di comprendere meglio il problema da risolvere, soprattutto grazie alle varie comunicazioni e incontri con il cliente. Al termine di questa fase sar\`a possibile quindi scrivere l'analisi dei requisiti che sar\`a la base per la fase di progettazione seguente. Queste attivit\`a potranno essere eseguite in modo parallelo su diverse parti del sistema, come anche le attivit\`a di progettazione e verifica, in modo da consentire la riduzione di tempi e costi di consegna del prodotto. Al termine delle attivit\`a di progettazione seguir\`a quella di programmazione, che produrr\`a il prototipo richiesto. L'ultima fase prima del collaudo sar\`a dedicata ad un attenta verifica di tutto il lavoro svolto nelle precedenti fasi. Il verificatore ed il responsabile di progetto sono le figure alle quali verranno affidate le responsabilit\`a pertinenti alle attivit\`a di verifica. Per una descrizione pi\`u dettagliata della pianificazione si veda il documento del PianoDiProgetto\_0\_5.pdf.

\section{Risorse necessarie, risorse disponibili}
Le attivit\`a di verifica necessitano di risorse umane e tecnologiche. Il gruppo \`e composto da sette membri, ognuno dei quali durante tutto il periodo di lavoro dovr\`a assumere, in periodi di tempo diversi (in base alla disponibilit\`a e soprattutto alle competenze di ciascuno) tutti i ruoli significativi per lo sviluppo del progetto. L'amministratore del progetto sar\`a tenuto a supervisionare tutte le fasi di verifica e a gestire le varie risorse necessarie per consentire un'attivit\`a di buon livello senza sprechi od oneri eccessivi. Per la comunicazione tra i componenti \`e stato creato un gruppo Google che permette la condivisione di files e la pianificazione di riunioni. Per l'archiviazione di grossi documenti o file verr\`a utilizzata la casella di posta elettronica del gruppo, mentre per i file di uso condiviso e quindi la cartella per la prossima revisione verr\`a usato lo spazio Google raggiungibile all'URL: \textit{http://groups.google.com/group/happycodeinc} e accessibile ai soli membri del gruppo.
Essendo inoltre ognuno dei membri provvisto di un account Gmail, verr\`a utilizzato il sistema di chat locale per la comunicazione interattiva.
\section{Strumenti, tecniche, metodi}
\subsection{Analisi Statica}
Questa fase verr\`a applicata durante la stesura del codice sorgente per rilevare incongruenze che possono sorgere tra il progetto e i requisiti. Verr\`a applicata inoltre sulla struttura delle varie componenti del sistema ''BR-jsys''. Essa comprende le seguenti sottofasi:
\begin{itemize}
\item \textbf{Analisi del flusso di controllo:} Verificher\`a una corretta esecuzione del codice  da parte delle varie componenti del prodotto.
In particolare, si controller\`a che:
\begin{itemize}
\item[-]il validatore accetti la regola in ingresso secondo le specifiche del linguaggio scelto;
\item[-]l'esecutore verifichi la consistenza dei dati espressi dagli oggetti business.
\end{itemize}
\item \textbf{Verifica formale del codice:} Verificher\`a la correttezza del codice scritto. Si constater\`a in particolare la correttezza totale di ogni modulo in modo che non conduca mai in uno stato di non terminazione.
\end{itemize}

\subsection{Analisi Dinamica}
Questa fase verr\`a applicata durante la progettazione e la stesura del codice al fine di verificare dinamicamente l'indipendenza delle singole unit\`a rispetto all'integrazione del sistema. Verr\`a testato quindi il sistema in tutti i suoi possibili casi e verranno effettuate prove per verificarne l'integrit\`a.
\begin{itemize}
\item Attraverso l'inserimento di nuove regole business in un contesto di prova, effettueremo test sulla validazione e li confronteremo con i risultati attesi.
\item Tramite opportuni oggetti business di prova effettueremo test sui risultati ottenuti dalla validazione, e li confronteremo con i risultati attesi.
\end{itemize}

\chapter[Gestione revisione]{Gestione amministrativa della revisione} 

\section{Comunicazione e risoluzione di anomalie}
Sar\`a compito dei verificatori la comunicazione delle anomalie e avverr\`a attraverso la stesura di documenti specifici. Costoro dovranno elencare le segnalazioni in modo che sia possibile individuare il punto errato in modo semplice e non ambiguo, eventualmente citando le parti non corrette. Discorso a parte va fatto per errori grammaticali o di formattazione che potranno essere corretti direttamente dai verificatori.

\section{Trattamento delle discrepanze} 
Se dovessero verificarsi discrepanze tra le necessit\`a del cliente ed i requisiti risultanti dall'analisi si provveder\`a ad un ulteriore analisi che avr\`a la priorit\`a sulle altre attivit\`a. Si provveder\`a all'aggiornamento dei documenti di analisi dei requisiti e se necessario ai prodotti delle attivit\`a successive all'analisi. L'analista, in collaborazione col progettista e l'amministratore, valuter\`a l'impatto delle modifiche sul lavoro gi\`a svolto e comunicher\`a ai membri interessati le variazioni.

\section{Procedure di controllo di qualit\`a di processo}
Il controllo della qualit\`a del software interessa l'intero processo di sviluppo del prodotto in questione. Esso comprende il monitoraggio e il miglioramento del processo, la verifica dell'applicazione di standard e procedure prestabiliti nonch\`e la soluzione degli eventuali problemi.
Per avere un esteso controllo degli errori, le modifiche apportate ai documenti, in seguito a inconsistenze riscontrate durante la fase di verifica e stesura, verranno notificate in ogni singolo cambiamento. Queste modifiche si troveranno all'inizio di ogni documento nella sezione ''Diario delle modifiche''. I verificatori sono tenuti a controllare parallelamente la documentazione sul lavoro svolto. Il responsabile di progetto deve verificare che le obiezioni sollevate dal verificatore siano corrette senza che si presentino discrepanze o anomalie. L'approvazione spetta pertanto al responsabile di progetto.

\chapter{Pianificazione ed esecuzione del collaudo}
\section{Specifica della campagna di validazione}
L'attivit\`a di collaudo \`e volta a qualificare il prodotto software sviluppato secondo i requisiti da noi identificati. Il risultato dell'attivit\`a di collaudo \`e determinante per procedere al rilascio del prodotto software. In base ai requisiti verranno applicate diverse tecniche di verifica:
\begin{itemize}
\item Alcuni dei requisiti tracciati dovranno essere validati mediante l'inserimento di un input e la verifica dell esattezza dei dati ottenuti. In particolare i requisiti 3.1.1.1, 3.1.1.2, 3.1.1.3, 3.1.1.4, 3.1.1.5, 3.1.1.6, 3.1.1.8, 3.1.2.1, 3.2.1.1, 3.2.1.4.
\item Altri invece necessitano di una approvazione da parte del committente attraverso incontri. In particolare i requisiti 3.2.1.2, 3.2.1.5, 3.2.1.6, 3.2.2.1, 3.2.2.2, 3.2.3.1, 3.2.3.2.
\item Infine vi sar\`a la necessit\`a di verificare la gestione corretta delle anomalie da parte del progetto. In particolare per il requisito 3.1.1.7.
\end{itemize}


\end{document}
	
